% 第14章 数値実装の設計（WP5 skeleton; 本文・導出は執筆予定）
\section{数値実装の設計}
\label{ch:ch14_numerics}

\paragraph{この章で導くこと（入力→出力）}
`cmbcore` の設計を本文として解説する（`03_numerics_spec.md` の人間可読版）。ソースを粗 k で解き k方向スプライン補間→細 k で視線積分する手法を含む。

% --- 節構成（01_textbook_spec.md より）-------------------------------------
\subsection{モジュール構成と単位系}
% TODO(WP5): 14.1 の本文。
\subsection{ODEソルバと tight coupling 切替}
% TODO(WP5): 14.2 の本文。
\subsection{視線積分の収束（粗解＋細k）}
% TODO(WP5): 14.3 の本文。
\subsection{検証プロトコルと CLASS 比較}
% TODO(WP5): 14.4 の本文。


\paragraph{導出契約}
本章の導出契約: —。各契約は「入力式→出力式」を中間ステップ省略なし
（1ステップ＝学部4年生が5分で検算可能）で履行する。\emph{本文プローズは WP5 で執筆。}

\paragraph{まとめ・チェックリスト・演習}
% TODO(WP5): 章末まとめ、チェックリスト、演習 3–6 問（うち1問は対応 NB の課題）。
